\honest{Living in Many Worlds}{Eliezer Yudkowsky}{2008}

Some readers are disturbed by the thought of constantly splitting into zillions of other people, which I call the ``straightforward and unavoidable prediction of quantum mechanics''. \nb{It is what the many-worlds reading says the theory describes. Single-world readings make the same predictions, and the book argues for many-worlds only later, in ``Many Worlds, One Best Guess''.} Others ask whether buckling a seat belt here makes another self unbuckle.

My answer is Egan's Law: ``It all adds up to normality.'' Frank Sulloway said that psychoanalysis beats Darwinism because its claims are outlandish enough to intrigue, while we feel we know Darwinism already, because we do. When Einstein replaced Newton's gravity, apples did not stop falling. Every new theory must reproduce what the old one got right, so many-worlds is not there to make strange predictions. Then why care? Because of the question: what adds up to normality? Quantum mechanics. If something else were there instead, the world would look strange. You have always lived in this universe.

Religions, anthropologists tell us, are minimally counterintuitive, surprising just enough to be memorable: Anubis has a dog's head on a man's body, and spirits see through walls but get hungry. Physics is not like that. Its underlying phenomena take long study, but its surface is ordinary. You will ``never'' glimpse another world or hear another self; that ``is unambiguously prohibited outright by the laws.'' ``Sorry, you're just schizophrenic.'' \nb{Five weeks earlier, in ``The Born Probabilities'', Yudkowsky had written that the interaction between branches ``may never actually become zero''. On the standard account, decoherence suppresses interference between worlds rather than forbidding it. The practical point stands.}

Decisions have nothing to do with branching. A decision feels like a branch point in your imagination, but you would feel the same uncertainty in one world, and people imagined alternatives long before quantum mechanics; a rock, which makes no decisions, splits too. You split all the time, not especially when you decide. In each world people choose what seems best to them, perhaps after different lines of thought; it is not that one version of you picks the best option and another the worst. You cannot choose which world to end up in. Choices in one world do not affect another, so buckling up here unbuckles no one elsewhere. I suggest treating worlds that are not in your future as a ``generalized past'', no more your responsibility than the twelfth century.

A 90\% chance and 9 worlds in 10 call for the same decision. \nb{That needs branch weights to act as probabilities. In ``The Born Probabilities'' Yudkowsky had said of that question: ``I don't know. It's an open problem.'' The attempt to answer it by Deutsch and Wallace is still contested.} Visualizing worlds may help people who struggle with probabilities, but it is ordinary decision theory, and a time to learn to shut up and multiply. The brain ``doesn't do 64-bit floating-point arithmetic,'' and between zero chance and epsilon chance ``there is an order-of-epsilon difference.'' A world with odds of a quadrillion to one probably exists if it is lawful, but deserves a quadrillionth of your feeling, or no thought at all; otherwise you might as well buy a lottery ticket with a quantum random number, a strategy guaranteed to produce a very tiny mega-win. To feel one in a trillion, roll a ten-sided die and think about the strange world only when it shows 9 twelve times running.

I allow a few implications for ethics. Average utilitarianism looks better, since ``there are already plenty of people exploring person-space.'' \nb{That persuades only someone who already thinks an added person is worth less when many exist. On the total view, other worlds change nothing.} And you can take joy in discovering anything you personally do not know, since being the ``first'' to know something means nothing when everything knowable is known in some other world. \nb{Yudkowsky had made that argument in ``Joy in Discovery''; it needs only a large universe.}

Otherwise, if your grasp of many-worlds is shaky and you are tempted by some strange belief, emotion or strategy: ``Don't.'' The quantum universe is the way things have always been.
